\section{A trace obstruction on the actual exceptional set}
Fix any finite real orders $q_k\ge1$ and any positive targets $H_k$.
There are half-noise chains in the two-base class, with exact matching
through $k$ and uniformly bounded stationary densities, for which the
operator $S$ of (F2) has norm tending to zero but
$\|S\|_{\mathcal S_{q_k}}\ge H_k$ already on the parity quotient.
Here the Schatten norms use the ordinary unnormalized trace. On that
quotient, even the exceptional multiplier has trivial stabilizer among
state permutations preserving the noise kernel. The dimensions may grow
arbitrarily rapidly. This rules out estimating the defect by requiring
such a finite trace norm to vanish uniformly in dimension.

Let $\ell$ be the least even integer greater than $k$ and $t=2^{-\ell}$.
Partition the ambient signs into $n+m$ disjoint groups of $\ell$ signs,
and keep their parities. There are $\ell$ disjoint bases. Every nonconstant
quotient Walsh character has ambient degree at least $\ell$, so any
normalized stationary density constructed from a quotient potential
matches through $k$. On the quotient, $P,Q$ are ordinary product noise
$T_t,T_{\sqrt t}$. Let $n=t^{-4}$, and call the first $n$ parities the
core, with the remaining $m$ coordinates free. For the core mean sign $d$
let $C=\{d\ge t/2\}$ and let $F$ consist of core configurations with at
most one negative sign, crossed with all free configurations. Set
$a=\exp(-nt^2/128)$. Take an unnormalized potential $h_0$ equal to one
on $C^c\cup F$ and $a$ on $C\setminus F$. Perturb each quotient value
upward by an amount strictly between zero and $a^2$, writing $H$ for
the result. Put $G=T_tH$, $Z=\mathbb E HG$, so $U=H/G$ and $\pi=HG/Z$.

The uniform probability of $C$ is at most $\exp(-nt^2/8)=a^{16}$.
From any starting point the probability of landing in $F$ is at most
$(n+1)((1+t)/2)^n\le a^9$ for large $k$. For a row in $F$ the mean
output core sign is at least $7t/8$ for large $n$.
The independent-sign Hoeffding bound gives probability at most $a^9$
of missing $C$. Thus on $F$, $a\le G\le a+a^2+2a^9$.
On $C^c$ the mean output sign is at most $t^2/2$ and the probability
of reaching $C$ is at most $\exp(-nt^2/32)=a^4$. There $G=1+O(a^2)$.
These estimates imply $Z=1+O(a^2)$ and $\pi\le3$ for large $k$.
On $C\setminus F$, $G\ge a$ gives $U\le1+a$; on $C^c$,
$U=1+O(a^2)$. On $F$ we have $(2a)^{-1}\le U\le2/a$.
Consequently $A=\{U>e^{\sigma k}\}=F$ for any fixed $\sigma>0$
and sufficiently large $k$.

In the quotient the $2$-to-$4$ inequality for $T_{\sqrt t}$ bounds
the positive sandwich norm by the $L^2$ norm of its multiplier.
The same bound holds for the full ambient norm by Perron invariance.
In particular
\[
 \|S\|\le (2/a)\sqrt{(n+1)2^{-n}}=o(1).
\]
The same estimate on $F$, the bound $O(1)$ for $|U-1|$ on
$C\setminus F$, and the uniform $O(a^2)$ bound on $C^c$ show
$\|QM_UQ-P\|=o(1)$. Thus the whole ambient spectrum recovers.

In positive semidefinite order on the quotient, $S$ dominates
$(2a)^{-1} T_{\sqrt t}^{(n)}M_{\{+\cdots+\}}T_{\sqrt t}^{(n)}
\otimes T_t^{(m)}$. The first sandwich has rank one and its
nonzero eigenvalue is the one-step return $((1+t)/2)^n$.
The eigenvalues of the second factor are $t^j$ with multiplicities
$\binom mj$. Monotonicity of each ordered eigenvalue yields
\[
 \|S\|_{\mathcal S_q}\ge (2a)^{-1}
       ((1+t)/2)^n(1+t^q)^{m/q}.
\]
For $q=q_k$ the last base exceeds one, so a finite $m$ achieves
any prescribed $H_k$. All preceding estimates are uniform in $m$.

Finally choose the perturbation within the open box so that every ratio
$H(y)/(T_tH)(y)$ is distinct. For any distinct $y,z$, equality gives
a nonzero polynomial condition: the coefficient of $H(y)^2$ in
$H(y)(T_tH)(z)-H(z)(T_tH)(y)$ is $T_t(z,y)>0$. A finite union of
proper polynomial zero sets cannot fill the box. Any symmetry of
$U{\bf1}_A$ therefore fixes each point of $F$. Since the noise entries
strictly distinguish Hamming distances, a noise-preserving permutation
is a Hamming isometry. The all-plus vertex and all of its individual
neighbors belong to $F$, and their distances distinguish every vertex.
The stabilizer is trivial.

The multiplicity factor is elementary tensor-product algebra. The
additional content is its realization on the \emph{actual} large-$U$
set of matched Doob chains, after quotienting out the ambient fibers,
with a vanishing killed norm and no residual permutation reduction.
Dimension-dependent order or a separate estimate that exploits
approximate rather than exact symmetry remains possible.
